% 【理论扩展】所属：第4章（由 chapters/revised/04-sources.tex \input 引入）
\minitopic{记忆保存论与生成论}保存论认为，记忆本身不创造最初确证，只把过去知识或确证传到现在。若主体过去不知道$p$，单纯记住同一内容不会变成知识。这解释记忆作为时间桥梁的作用，也防止反复回想任意信念提高地位。 生成论指出，记忆可重新组织材料、保留内容而丢失过去的缺陷，或与其他信息结合产生新知识\parencite{lackey2005}。例如儿童最初因不可靠传闻相信历史日期，成年后忘记来源，却在可靠教材中再次见到并形成稳定记忆；当前知识不应由最初缺陷永久污染。需识别真正维持当前信念的链条。 记忆还可能提供关于过去经验的证据，即使过去未形成显式信念。目击者当时注意到一件蓝外套却未判断身份，后来回忆细节并结合照片作出推断。这里记忆保存经验信息，当前推理生成命题知识。

\minitopic{记忆内容与真实性}情景记忆常包含“再经历”感，语义记忆则保存事实而不带具体场景。两者错误模式不同。鲜明情景感会增加确信，却不保证准确；语义记忆可长期可靠，却失去来源标签。认识规范应针对记忆类型。 回忆不是从固定仓库取出副本。当前目标、后来叙述和提问方式会参与重构；误导信息效应已有长期实验研究支持\parencite{loftus2005}。重构性并不推出记忆普遍不可靠；知觉同样经过加工。关键是哪些机制在何种条件下系统偏离，以及是否有外部校正。 “记忆一致”也不等于独立证实。多人讨论后形成相似叙述，可能共享污染。调查应尽早分别记录，保存首次陈述与修改轨迹。制度设计将心理学事实转化为认识保障。

\minitopic{遗忘来源与认识信用}日常知识大量处于来源遗忘状态。若每次都要求回忆教师、书名与页码，认知负担不可承受。默认可依赖稳定记忆，直到出现具体击败者；高争议事实、精确数字和直接引语则需重新查源。 来源监控能力可分级。主体至少区分亲历、听闻和推断；更高要求是识别具体来源可靠性。数字环境把这项能力外部化：浏览历史、文献管理器和版本记录可作为扩展记忆。工具失效或记录被篡改也会成为新风险。 认识信用可以分布。研究者记得结论但依赖实验室笔记恢复理由，知识不必完全存于头脑。若记录系统可追溯、主体知道如何访问并能响应质疑，它可能构成记忆实践的一部分。

\minitopic{内省的方法模型}内省可被理解为内部观察、直接觉知、理性自我归属或解释推断。内部观察模型把心理状态当作内在对象，却面临需要另一个“内眼”的回溯。直接觉知避免推理，却较难解释错误。理性模型认为，回答“我相信$p$吗”通常通过考虑$p$本身，而非搜寻内部迹象\parencite{moran2001}。 透明性方法适合当前判断，却不覆盖压抑动机、性格倾向和无意识偏见。对于这些对象，我们常像研究他人一样，根据行为、情境和反馈推断。自知不是单一通道，权威随对象变化。 第一人称判断还有表达与构成作用。说“我决定支持方案”可部分形成承诺，而说“我其实嫉妒”更像发现解释。把所有自我报告当作描述事实，会误解意志与承诺；把所有都当作不可错宣告，则掩盖自欺。

\minitopic{自我欺骗与认识责任}自我欺骗案例中，主体选择性注意有利证据、提高不利证据门槛，却维持“我只是客观”。信念不完全受直接意志控制，但信息搜索和解释习惯可受间接控制。责任来自主体是否有机会认识偏差并建立纠正程序。 并非所有错误自我评价都是自欺。能力有限、情绪识别困难或社会词汇缺失都可能造成无责错误。责备前需说明动机性偏差如何影响证据处理，而不能从“错误恰好有利于他”直接推断故意。 写反方备忘录、预先规定证据门槛、请利益独立者评估，都是针对自我欺骗的外部支架。认识自主包含主动限制自己的偏差机会。

\minitopic{分析性与先验性}分析命题通常被理解为凭意义为真，如“所有单身者未婚”；先验知识则按确证来源分类。若一个分析真理因权威证言而被相信，这次信念的实际根据可能是经验性证言；主体原则上可先验知道，不表示每个主体的每次知识都先验。 Quine等人质疑分析与综合之间有清楚边界。词义信念与世界信念处在同一可修正网络，面对足够压力甚至逻辑规则也可调整。回应者可承认边界不总锋利，却保留典型案例中不同辩护方式。 概念能力本身通过经验获得，不妨碍使用后进行先验推理。学习“如果”和“所有”需要语言互动；一旦掌握，从“所有A是B”和“x是A”推出“x是B”的理由不依赖再次观察世界。学习因果史与当前确证基础必须分开。

\minitopic{理性直觉与其可靠性}理性主义常诉诸对必然关系的理性直觉\parencite{bonjour1998}。反对者问，直觉为何可靠接触抽象真理？若只因“显得必然”便可信，与知觉“显得如此”有何相同和差异？一种答案把理性能力视为理解概念时的构成性把握，另一种以最佳解释或演化能力说明。 直觉也会冲突。集合论悖论、概率错误和复杂思想实验显示，初步显现可被证明击败。温和理性主义把直觉视为可错的初步理由，需与证明、理论统一和他人批评结合，而非神谕。 形式证明提供可公开检查的步骤，但对公理、规则和符号理解仍有背景依赖。先验知识不是完全私人；共享证明实践和纠错制度同样重要。

\minitopic{归纳问题的不同回答}休谟指出，从过去规律到未来规律既非演绎有效，也不能用过去归纳成功无循环证明\parencite{hume2000}。一种回应是实用辩护：若世界有规律，归纳最终有机会发现；若无规律，任何方法都无能。它说明使用归纳的策略优势，却未必证明具体预测为真。 规则循环主义认为，一种基本推理规则可用自身支持而不恶性，只要没有不使用该规则的独立可行辩护。外在主义则说，归纳实际在相对稳定世界可靠，主体不需先证明总体可靠。贝叶斯方法把问题转为先验与更新，但仍假设条件化和模型结构。 自然规律与因果知识可增强具体归纳。不是只因太阳过去升起便预测，而是借助天体模型；不是只看药物过去有效，而是理解机制和试验设计。解释结构不能彻底消除归纳，因为模型本身也由经验支持，但能区分强弱推断。

\minitopic{综合案例：家族叙事}一个人鲜明记得祖母说家族在某年迁居。多年重复讲述使记忆稳定，地方档案却给出不同年份。分析要区分：他记得祖母曾说，还是迁居事实本身；重复来自独立回忆还是每次重述；档案是否完整；年份差异是否由历法或登记延迟解释。 他可以知道“祖母讲过这个故事”，却不知道故事内容为真。若后来找到同期信件，证言、记忆和文档汇合；若信件正是祖母根据同一错误记忆写成，来源仍不独立。认识链的谱系比材料数量重要。 比较视角还提醒，家族记忆承担身份和伦理功能。质疑事实准确性不必否定叙事价值；尊重叙事也不应把象征意义当作历史证据。区分功能使讨论既不粗暴拆穿，也不放弃真值评价。

\begin{tcolorbox}[breakable,colback=white,colframe=blue!30!black,title={记忆与先验分析表},fonttitle=\bfseries]
记忆部分：过去是否知道；保存的是内容、经验还是来源；当前信念由何链条维持；有无重构和社会污染。先验部分：区分学习条件与确证基础、分析与必然、直觉与证明；写出可能击败者及公开校正方式。
\end{tcolorbox}
